%
%Academic CV LaTeX Template
% Author: Dubasi Pavan Kumar
% LinkedIn: https://www.linkedin.com/in/im-pavankumar/
% License: MIT
%
% For errors, suggestions, or improvements, please contact:
% Email: pavankumard.pg19.ma@nitp.ac.in
%



\documentclass[a4paper,11pt]{article}

% Package imports
\usepackage{latexsym}
\usepackage{xcolor}
\usepackage{float}
\usepackage{ragged2e}
\usepackage[empty]{fullpage}
\usepackage{wrapfig}
\usepackage{lipsum}
\usepackage{tabularx}
\usepackage{titlesec}
\usepackage{geometry}
\usepackage{marvosym}
\usepackage{verbatim}
\usepackage{enumitem}
\usepackage{fancyhdr}
\usepackage{multicol}
\usepackage{graphicx}
\usepackage{DejaVuSerif}
\usepackage{DejaVuSans}
\usepackage{DejaVuSansMono}
%\usepackage{cfr-lm}
\usepackage[T1]{fontenc}
\usepackage{fontawesome5}

% Color definitions
\definecolor{darkblue}{RGB}{0,0,139}

% Page layout
\setlength{\multicolsep}{0pt} 
\pagestyle{fancy}
\fancyhf{} % clear all header and footer fields
\fancyfoot{}
\renewcommand{\headrulewidth}{0pt}
\renewcommand{\footrulewidth}{0pt}
\geometry{left=1.4cm, top=0.8cm, right=1.2cm, bottom=1cm}
\setlength{\footskip}{5pt} % Addressing fancyhdr warning

% Hyperlink setup (moved after fancyhdr to address warning)
\usepackage{xurl}
\usepackage[hidelinks]{hyperref}
\hypersetup{
    colorlinks=true,
    linkcolor=darkblue,
    filecolor=darkblue,
    urlcolor=darkblue,
}

% Custom box settings
\usepackage[most]{tcolorbox}
\tcbset{
    frame code={},
    center title,
    left=0pt,
    right=0pt,
    top=0pt,
    bottom=0pt,
    colback=gray!20,
    colframe=white,
    width=\dimexpr\textwidth\relax,
    enlarge left by=-2mm,
    boxsep=4pt,
    arc=0pt,outer arc=0pt,
}

% URL style
\urlstyle{tt}
\makeatletter
\let\HREF\href
\renewcommand{\href}[2]{\HREF{#1}{\begingroup\UrlFont #2\endgroup}}
\makeatother
% Text alignment
\raggedright
\setlength{\tabcolsep}{0in}

% Section formatting
\titleformat{\section}{
  \vspace{-4pt}\scshape\raggedright\large
}{}{0em}{}[\color{black}\titlerule \vspace{-7pt}]

% Custom commands
\newcommand{\resumeItem}[2]{
  \item{
    \textbf{#1}{\hspace{0.5mm}#2 \vspace{-0.5mm}}
  }
}

\newcommand{\resumePOR}[3]{
\vspace{0.5mm}\item
    \begin{tabular*}{0.97\textwidth}[t]{l@{\extracolsep{\fill}}r}
        \textbf{#1}\hspace{0.3mm}#2 & \textit{\small{#3}} 
    \end{tabular*}
    \vspace{-2mm}
}

\newcommand{\resumeSubheading}[4]{
\vspace{0.5mm}\item
    \begin{tabular*}{0.98\textwidth}[t]{l@{\extracolsep{\fill}}r}
        \textbf{#1} & \textit{\footnotesize{#4}} \\
        \textit{\footnotesize{#3}} &  \footnotesize{#2}\\
    \end{tabular*}
    \vspace{-2.4mm}
}

\newcommand{\resumeProject}[4]{
\vspace{0.5mm}\item
    \begin{tabular*}{0.98\textwidth}[t]{l@{\extracolsep{\fill}}r}
        \textbf{#1} & \textit{\footnotesize{#3}} \\
        \footnotesize{\textit{#2}} & \footnotesize{#4}
    \end{tabular*}
    \vspace{-2.4mm}
}

\newcommand{\resumeSubItem}[2]{\resumeItem{#1}{#2}\vspace{-4pt}}

\renewcommand{\labelitemi}{$\vcenter{\hbox{\tiny$\bullet$}}$}
\renewcommand{\labelitemii}{$\vcenter{\hbox{\tiny$\circ$}}$}

\newcommand{\resumeSubHeadingListStart}{\begin{itemize}[leftmargin=*,labelsep=1mm]}
\newcommand{\resumeHeadingSkillStart}{\begin{itemize}[leftmargin=*,itemsep=1.7mm, rightmargin=2ex]}
\newcommand{\resumeItemListStart}{\begin{itemize}[leftmargin=*,labelsep=1mm,itemsep=0.5mm]}

\newcommand{\resumeSubHeadingListEnd}{\end{itemize}\vspace{2mm}}
\newcommand{\resumeHeadingSkillEnd}{\end{itemize}\vspace{-2mm}}
\newcommand{\resumeItemListEnd}{\end{itemize}\vspace{-2mm}}
\newcommand{\cvsection}[1]{%
\vspace{2mm}
\begin{tcolorbox}
    \textbf{\large #1}
\end{tcolorbox}
    \vspace{-4mm}
}

\newcolumntype{L}{>{\raggedright\arraybackslash}X}%
\newcolumntype{R}{>{\raggedleft\arraybackslash}X}%
\newcolumntype{C}{>{\centering\arraybackslash}X}%

% Commands for icon sizing and positioning
\newcommand{\socialicon}[1]{\raisebox{-0.05em}{\resizebox{!}{1em}{#1}}}
\newcommand{\ieeeicon}[1]{\raisebox{-0.3em}{\resizebox{!}{1.3em}{#1}}}

\begin{document}

% Header
\begin{center}
	{\large\textbf{Benjamin T. James}}
\end{center}
\vspace{-6mm}

\begin{center}
	\small{
		\socialicon{\faPhone} (314)-703-5233 | \socialicon{\faEnvelope} \href{mailto:benjames@mit.edu}{benjames@mit.edu} |
		\socialicon{\faUser} \url{https://www.mit.edu/~benjames/}
	}
\end{center}
\vspace{-6mm}
\begin{center}
	\small{
		\socialicon{\faBookReader} \href{https://scholar.google.com/citations?user=t0y3zRkAAAAJ}{Google Scholar}
		\socialicon{\faGithub} \href{https://github.com/benjamin-james}{benjamin-james} |
		\socialicon{\faTwitter} \href{https://twitter.com/\_\_ben\_james\_\_}{@\_\_ben\_james\_\_}
	}
\end{center}

\vspace{-4mm}

%% \section{\textbf{Overview}}
%% \vspace{1mm}
%% \small{
%%   Ph.D student seeking a challenging position to leverage my expertise in computational genomics.
%%   Aiming to contribute to innovative projects at the intersection of AI and practical problem-solving in fields such as [specific areas of interest].
%% }
%% \vspace{-2mm}



%% \section{\textbf{Experience}}
%% \vspace{-0.4mm}
%%   \resumeSubHeadingListStart
%%   \resumeSubheading
%%       {{Company A [\href{https://www.companya.com}{\faIcon{globe}}]}}{City, Country}
%%       {Job Title A}{Month Year - Month Year}
%%       \resumeItemListStart
%%         \item Developed [specific achievement] achieving [specific metric] in [specific area]
%%         \item Implemented [technology/method], enhancing [specific aspect] by [specific percentage]
%%         \item Conducted analysis on [specific data], identifying [key findings]
%%         \item Presented findings at [specific event], receiving [specific recognition]
%%       \resumeItemListEnd 
%%   \resumeSubheading
%%     {Company B [\href{https://www.companyb.com}{\faIcon{globe}}]}{Remote}
%%     {Job Title B}{Month Year - Month Year}
%%     \resumeItemListStart
%%       \item Engineered a [specific system/model], improving [specific metric] by [percentage]
%%       \item Developed [specific tool/method], increasing [specific aspect] by [percentage]
%%       \item Implemented [specific system], reducing [specific metric] by [percentage]
%%       \item Conducted [specific test/analysis] to validate [specific aspect]
%%     \resumeItemListEnd
%%   \resumeSubHeadingListEnd
%% \vspace{-6mm}

\section{\textbf{Education}}
\vspace{-0.4mm}
\resumeSubHeadingListStart
\resumeSubheading
{Massachusetts Institute of Technology}{Cambridge, Massachusetts}
{Ph.D in Electrical Engineering and Computer Science, minor in Econometrics}{September 2019 -- August 2026}
\resumeItemListStart
\item GPA: 4.80/5.00
\resumeItemListEnd

\resumeSubheading
{University of Tulsa}{Tulsa, Oklahoma}
{BSc in Mathematics, BSc in Computer Science}{Aug 2015 -- May 2019}
\resumeItemListStart
\item GPA: 3.747/4.000
\resumeItemListEnd

%% \resumeSubheading
%% {}{City, Country}
%% {Pre-University Education}{Month Year}
%% \resumeItemListStart
%% \item Grade: XX.X\%
%% \resumeItemListEnd

%% \resumeSubheading
%% {High School Name}{City, Country}
%% {Secondary Education}{Month Year}
%% \resumeItemListStart
%% \item GPA: X.X/10
%% \resumeItemListEnd

\resumeSubHeadingListEnd
\vspace{-6mm}

\section{\textbf{Selected Research} \small{Includes two (co-) first author \textbf{Cell} papers, one second-author \textbf{Nature} paper}}
\vspace{-0.4mm}
\resumeSubHeadingListStart

\resumeProject
{Single-cell dissection of the human brain}
{Tools: R, Python, Bash, C++}
{2021 - 2026}
{{}[\href{https://github.com/kellislab/}{\textcolor{darkblue}{\faGithub}}]}
\resumeItemListStart
\item Integrated high-dimensional, multimodal (scRNA-seq/scATAC-seq) data from \textbf{hundreds of postmortem human brains} and mouse brains from corresponding model systems
\item Analyzed these data across several \textbf{human disease contexts} and brain regions using \textbf{functional genomics}, careful integration/atlasing, and case-control causal inference to infer disease pathology.  %\textcolor{red}{TODO highlight keywords}
\item Created new methods to make inferences on new technologies (e.g. spatial transcriptomics) and gene module discovery, mostly using random matrix theory and high-dimensional statistics
\item Connected millions of DNA elements active across human tissues to their target genes

\item \tiny{Linville*, R. M., \textbf{James*, B. T.}, Galani, K., Ho, L., Shin, J., Oliver, E., Bock, R., Murray, E. M., Louca, M., Oke, O., Wang, C., Engelberg-Cook, E., DeTure, M., Farrell, V., Fitzwalter, B., Pineda, S., Sathitloetsakun, S., Liang, X., Madras, B., and Dickson, D. W., Mash, D. C., Turecki, G., Wheeler, V. C., Alvarez, V. A., Gabuzda, D., Kellis, M., Heiman, M. (2025). Cross-species single cell atlas of the striatum defines cell type and subregion disease vulnerabilities. \textbf{Cell}}.
\item \tiny{Xiong*, X., \textbf{James*, B. T.}, Boix*, C. A., Park, Y. P., Galani, K., Victor, M. B., Sun, N., Hou, L., Ho, L.-L., Mantero, J., Scannail, A. N., Dileep, V., Dong, W., Mathys, H., Bennett, D. A., Tsai, L.-H., \& Kellis, M. (2023). Epigenomic dissection of Alzheimer’s disease pinpoints causal variants and reveals epigenome erosion. \textit{Cell}, \textit{186}(20), 4422–4437.e21. \url{https://doi.org/10.1016/j.cell.2023.08.040}}
\item \tiny{Boix, C. A., \textbf{James, B. T.}, Park, Y. P., Meuleman, W., \& Kellis, M. (2021). Regulatory genomic circuitry of human disease loci by integrative epigenomics. \textit{Nature}, \textit{590}(7845), 300–307. \url{https://doi.org/10.1038/s41586-020-03145-z}}
\resumeItemListEnd

%% \resumeProject
%% {Scalable, spatially-informed clustering for spatial transcriptomics}
%% {Tools: Python, ARPACK (scipy)}
%% {2025-2026}
%% {{}[\href{https://github.com/kellislab/banksy-hd}{\textcolor{darkblue}{\faGithub}}]}
%% \resumeItemListStart
%% \item Adaptation of the \href{https://doi.org/10.1038/s41588-024-01664-3}{BANKSY algorithm} for high dimensional transcriptomics
%% \item Newer spatial transcriptomics assays like 10x Genomics Visium HD contain millions of spots per experiment that are not amenable to these methods
%% \item Tech specific adaptations include a Sinkhorn-Knopp based noise removal (de-striping) and a Laplacian of Gaussian filter (for blob detection) replacing a Gabor filter
%% \item Sparse incremental Lanczos-based PCA that allows for input matrices to stay sparse, as opposed to the sklearn based IncrementalPCA
%% \item Designed and implemented alone, with no collaborators and no advisor input
%% \resumeItemListEnd

%% \resumeProject
%% {Linking regulatory elements to their target genes}
%% {Tools: Python, R, XGBoost}
%% {2019-2023}
%% {{}[\href{https://github.com/kellislab/epimap}{\textcolor{darkblue}{\faGithub}}]}
%% \resumeItemListStart
%% \item Developed a method that links annotated enhancers to target genes by using correlated activity between epigenomes (imputed ChIP-seq) over putative enhancers and corresponding gene expression \textbf{[J.12]},\textbf{[J.9]},\textbf{[P.18]},\textbf{[J.2]}
%% \item Started as rotation project in Manolis Kellis lab, directly supervised under Carles Boix (1st author of Nature paper), where my contributions placed me as a second author on a Nature paper \textbf{[J.12]}
%% \item Worked with the ENCODE Distal Regulation group and GENCODE with this method to annotate regulatory elements \textbf{[J.8]},\textbf{[J.4]},\textbf{[P.19]}
%% \item Later this methodology for linking putative enhancers to genes was extended for single-nucleus ATAC-seq studies for Alzheimer's disease \textbf{[J.11]},\textbf{[J.10]}
%% \resumeItemListEnd

\resumeProject
{Intelligent clustering of DNA sequences}
{Tools: C++/OpenMP; gdb/valgrind}
{2016-2019}
{{}[\href{https://github.com/BioinformaticsToolsmith/MeShClust}{\textcolor{darkblue}{\faGithub}}]}
\resumeItemListStart
\item Developed a fast, concurrent method as an undergraduate student to approximate an expensive method for DNA similarity (alignment) and used for clustering. Led to \textbf{5} journal publications as an undergrad.
\item \tiny{\textbf{James, B. T.}, Luczak, B. B., \& Girgis, H. Z. (2018). MeShClust: an intelligent tool for clustering DNA sequences. \textit{Nucleic Acids Research}, \textit{46}(14), e83. \url{https://doi.org/10.1093/nar/gky315}}
\resumeItemListEnd

\resumeSubHeadingListEnd

\section{\textbf{Skills}}
\vspace{-0.4mm}
\resumeHeadingSkillStart
\resumeSubItem{Programming Languages:}
{R and Python (Daily use), C++ (+pybind11/Rcpp), Golang, Bash/POSIX sh, C, AWK \& Perl, Julia, Common Lisp}
\resumeSubItem{Research-specific Packages and Tools:}
{Bioconductor, scverse, samtools/bedtools, Eigen, NumPy/SciPy/SK-Learn, Snakemake \& Nextflow, brms}
\resumeSubItem{DevOps/MLOps:}
{De-facto manager of two lab-wide compute clusters, managing software (Lmod, conda, Docker, singularity environments, CUDA/NV admin, Terraform, Ansible, etc). Experienced Linux user (Debian/Gentoo user since 2014), see \href{https://github.com/benjamin-james/dotfiles}{my dotfiles}, formerly an ArchLinux \href{https://aur.archlinux.org/packages?SeB=m&K=benjamin-james}{maintainer}}
\resumeSubItem{Side projects:}
{Wrote an \textbf{agentic} harness tool to specify disposable agent virtual machines \href{https://github.com/benjamin-james/agentctl}{\faGithub} and deploy them \href{https://github.com/benjamin-james/dotfiles/blob/master/bin/vm}{\faGithub},
	Contributed the dot-product distance into the \href{https://github.com/eddelbuettel/rcppannoy}{RcppAnnoy} package required by R genomics libraries \href{https://github.com/eddelbuettel/rcppannoy/pull/78}{\faGithub},
	Wrote the sparse-matrix math for the \href{https://github.com/pinellolab/pychromVAR}{pychromVAR} package for transcription factor enrichment in snATAC-seq \href{https://github.com/pinellolab/pychromVAR/pull/13}{\faGithub},
	wrote a complete BF compiler \href{https://github.com/benjamin-james/brainfuck}{\faGithub} and a Pascal compiler frontend \href{https://github.com/benjamin-james/pascal-compiler}{\faGithub} both in C, wrote and deployed a multi-user webapp on a VPS \href{https://github.com/benjamin-james/cfb_survivor_pool}{\faGithub} (Apache/Flask/PostgresQL)}
\resumeHeadingSkillEnd

\section{\textbf{Honors and Awards}}
\vspace{-0.4mm}
\resumeSubHeadingListStart

\resumeProject
{NSF GRFP}
{National Science Foundation}
{May 2019}
{{}}%[\href{https://award-link-a.com}{\textcolor{darkblue}{\faIcon{globe}}}]}

\resumeProject
{Cronin Fellowship}
{MIT EECS department}
{April 2019}
{{}}
%% \resumeItemListStart
%% \item MIT departmental (EECS) fellowship for exceptional Ph.D applicants
%% \resumeItemListEnd
\resumeSubHeadingListEnd

\vspace{-6mm}
%% \section{\textbf{Leadership Experience}}
%% \vspace{-0.4mm}
%% \resumeSubHeadingListStart
%% \resumeProject
%%   {Leadership Role A}
%%   {Organization/Institution Name}
%%   {Month Year - Month Year}
%%   {{}[\href{https://organization-a-link.com}{\textcolor{darkblue}{\faIcon{globe}}}]}
%% \resumeItemListStart
%%   \item Key responsibility or achievement in this role
%%   \item Quantifiable impact or improvement made during tenure
%%   \item Initiative taken or project led
%% \resumeItemListEnd

%% \resumeProject
%%   {Leadership Role B}
%%   {Organization/Institution Name}
%%   {Month Year - Month Year}
%%   {{}[\href{https://organization-b-link.com}{\textcolor{darkblue}{\faIcon{globe}}}]}
%% \resumeItemListStart
%%   \item Key responsibility or achievement in this role
%%   \item Quantifiable impact or improvement made during tenure
%%   \item Initiative taken or project led
%% \resumeItemListEnd

%% \resumeSubHeadingListEnd

%% \vspace{-6mm}

%% \section{\textbf{Volunteer Experience}}
%% \vspace{-0.4mm}
%% \resumeSubHeadingListStart
%% \resumeProject
%%   {Volunteer Role A}
%%   {Organization Name}
%%   {Month Year - Month Year}
%%   {{}[\href{https://volunteer-org-a-link.com}{\textcolor{darkblue}{\faIcon{globe}}}]}
%% \resumeItemListStart
%%   \item Key responsibility or contribution in this role
%%   \item Impact of your volunteer work
%%   \item Skills developed or applied during this experience
%% \resumeItemListEnd

%% \resumeProject
%%   {Volunteer Role B}
%%   {Organization Name}
%%   {Month Year - Present}
%%   {{}[\href{https://volunteer-org-b-link.com}{\textcolor{darkblue}{\faIcon{globe}}}]}
%% \resumeItemListStart
%%   \item Key responsibility or contribution in this role
%%   \item Impact of your volunteer work
%%   \item Skills developed or applied during this experience
%% \resumeItemListEnd

%% \resumeSubHeadingListEnd
%% \vspace{-6mm}

%% \section{\textbf{Professional Memberships}}
%% \vspace{-0.4mm}
%% \resumeSubHeadingListStart
%% \resumePOR{Professional Organization A}
%%     {, Membership ID: XXXXXXXX}
%%     {Month Year - Present}
%% \resumePOR{Professional Organization B}
%%     {, \href{https://membership-certificate-link.com}{Membership ID: XXXXXXXX}}
%%     {Month Year - Present}
%% \resumePOR{Professional Organization C}
%%     {, \href{https://membership-certificate-link.com}{Membership ID: XXXXXXXX}}
%%     {Month Year - Present}

%% \resumeSubHeadingListEnd
%% \vspace{-6mm}

%% \section{\textbf{Certifications}}
%% \vspace{-0.2mm}
%% \resumeSubHeadingListStart
%% \resumePOR{}{\href{https://certification-link-a.com}{
%% \textbf{Certification A}
%% }}{Month Year}
%% \resumePOR{}{
%% \textbf{Certifying Body:} {{\href{https://certification-link-b.com}{Certification B}}}}{Month Year}
%% \resumePOR{}{
%% \textbf{Certifying Body:} {{\href{https://certification-link-c.com}{Certification C}}}}{Month Year}
%% \resumePOR{}{\href{https://certification-link-d.com}{
%% \textbf{Certification D}
%% }}{Month Year}

%% \resumeSubHeadingListEnd
%% \vspace{-6mm}

%% \section{\textbf{Additional Information}}
%% \vspace{-0.4mm}
%% \small{
%% \textbf{Languages:} Language A (Proficiency level), Language B (Proficiency level), Language C (Proficiency level)

%% \textbf{Interests:} Interest 1, Interest 2, Interest 3, Interest 4
%% }
%% \vspace{-4mm}


%% \section{\textbf{References}}
%% \vspace{-0.2mm}
%% \small{
%% \begin{enumerate}[leftmargin=*,labelsep=2mm]
%% \item \textbf{Reference Person 1}\\
%%    Job Title, Department\\
%%    Organization/Institution Name\\
%%    Email: email1@example.com\\
%%    Phone: +X-XXX-XXX-XXXX\\
%%    \textit{Relationship: [e.g., Thesis Advisor, Manager, etc.]}

%% \item \textbf{Reference Person 2}\\
%%    Job Title, Department\\
%%    Organization/Institution Name\\
%%    Email: email2@example.com\\
%%    Phone: +X-XXX-XXX-XXXX\\
%%    \textit{Relationship: [e.g., Project Supervisor, Colleague, etc.]}

%% \item \textbf{Reference Person 3}\\
%%    Job Title, Department\\
%%    Organization/Institution Name\\
%%    Email: email3@example.com\\
%%    Phone: +X-XXX-XXX-XXXX\\
%%    \textit{Relationship: [e.g., Mentor, Collaborator, etc.]}
%% \end{enumerate}
%% }

\end{document}
